\head{\textsc{\emph{Ilíada}, Canto III}}
\section*{\centering\textbf{\textsc{Canto} III}\footnote{\ehuno, pp. 43-46, 49, 52-53, 62 + \ehdos, p. 81 + \neh, p. 24.}}
\addcontentsline{toc}{subsection}{\textsc{Canto III}}
%
\noindent-- -- -- -- -- --\par
%
\beginnumbering
\pstart
\setline{21}
Y el marcial Menelao, cuando lo vé que al frente\\
de las tropas, avanza con andar imponente,\\
se alegra como hambriento león que hace gran presa\\
de un enastado ciervo o una cabra montesa,\\
y ávido los devora, ni aunque esté perseguido\\
por los ágiles perros y los mozos fornidos.\par
\pend
%
\noindent-- -- -- -- -- --\par
%
\pstart
\setline{125}
Hallóla en el palacio tejiendo una amplia y bella\\
tela doble de púrpura en que historiaba entonces,\\
las mil luchas que a influjo de Ares sufren por ella\\
domadores troyanos y aqueos que arma el bronce.\par
\pend
%
\noindent-- -- -- -- -- --\par
%
\pstart
\setline{136}
El marcial Menelao y Alejandro, al instante,\\
por ti van a batirse con sus enormes lanzas,\\
y tú serás la esposa del que salga triunfante.\\
Diciendo así la diosa, con dulces esperanzas\\
de su primer marido, patria y deudos, la arroba.\\
Y al punto, envuelta en cándidos velos, deja la alcoba,\\
vertiendo tierno llanto...\par
\pend
%
\noindent-- -- -- -- -- --\par
%
\pstart
\setline{154}
Al ver Helena hacia ellos marchaba, con voz queda\\
y en aladas palabras dijeron: --Justo es para\\
los troyanos y aqueos padecer tantos males\\
por tal mujer; pues a las diosas inmortales,\\
ella, terriblemente, se parece de cara.\par
\pend
%
\noindent-- -- -- -- -- --\par
%
\pstart
\setline{164}
Tú no eres la culpable, sino los dioses, de esta\\
guerra de los aqueos, de plorable y funesta.\par
\pend
%
\noindent-- -- -- -- -- --\par
%
\pstart
\setline{173}
Dulce me habría sido la peor muerte, cuando\\
por seguir aquí a tu hijo, dejé el tálamo amable,\\
hermanos, dulce hija y amigas adorables.\\
Y de que así no fuera, consúmome llorando.\par
\pend
%
\noindent-- -- -- -- -- --\par
%
\pstart
\setline{178}
Es el poderosísimo Agamenón Atrida,\\
a la vez tan cumplido rey como buen guerrero:\\
cuñado de esta perra, por mi honra inmerecida.\par
\pend
%
\noindent-- -- -- -- -- --\par
%
\pstart
\setline{241}
\firstlinenum{242}
¿O no querrán, acaso, concurrir al combate,\\
corridos del oprobio que sobre mí se abate?\par
\pend
%
\noindent-- -- -- -- -- --\par
%
\pstart
\setline{320}
\firstlinenum{320}
Zeus padre que en el Ida reinas glorioso y grande:\\
haz que quien puso entre ambos pueblos el desconcierto,\\
descienda al fondo de la morada de Hades, muerto;\\
y que amistad, en cambio, tu alianza fiel nos mande.\par
\pend
%
\noindent-- -- -- -- -- --\par
%
\pstart
\setline{428}
Huiste de la guerra donde habría debido\\
inmolarte aquel héroe que antes fué mi marido.\\
A Menelao, amado de Ares, te lisonjeabas\\
de exceder en corage, brazo y lanza. Anda ahora\\
y rétalo de nuevo, como cuando peleabas...\\
Mas, nó, que disuadirte querría, previsora,\\
para que locamente no afrontes la pujanza\\
del blondo Menelao; no sea que éste al punto,\\
en la contienda bélica te rinda por la lanza.\par
\pend
%
\noindent-- -- -- -- -- --\par
%
\pstart
\setline{438}``Mujer, no me zahieras así. Si venció ahora\\
Menelao, con ayuda de Atena, vendrá mi hora;\\
porque también hay dioses que a nosotros nos guían.\\
Ea, pues, acostémonos y a querernos volvamos,\\
pues nunca embargó tanto tu amor el alma mía,\\
ni cuando de la amena Lacedemonia un día\\
te rapté, y en las naves, pasando el mar, bogamos,\\
y afecto y lecho uniéronnos en la lista de Kranae.\\
Así te amo y el ducle deseo a ti me atrae.\\
Tal dijo, y fuese al lecho, seguido de la esposa,\\
y ambos allá acostáronse en la cama lujosa''.
\pend
%
\noindent-- -- -- -- -- --\par
%
\endnumbering